\section*{第 \ref{ch:b2:structures} 章 --- 集合与结构}

\begin{solution}{exo:b2:structures:1}
\emph{$\N$ 的有限子集：}可数 --- $\intint{0}{n}$ 的子集全体
为有限集，有限子集恰为这些集合的可数并
对 $n$ 取并
（\cref{prop:b2:structures:countablestable}~(3)）；由于
它包含全部单点集，故无限。

\emph{$\N$ 的全体子集：}不可数，由康托定理
（\cref{thm:b2:structures:cantor}~(1)，取 $E = \N$）。

\emph{$\R \setminus \Q$：}不可数 --- 否则 $\R = \Q \cup
(\R\setminus\Q)$ 将是两个可数集之并，
与 \cref{thm:b2:structures:cantor}~(2) 矛盾。

\emph{$\Q$ 上的多项式：}可数 --- 次数
$\leq n$ 的多项式可单射入 $\Q^{n+1}$（有限个可数集的积），
再对 $n$ 取并。

\emph{最终为零的 0-1 数列：}可数 --- 它们与
$\N$ 的有限子集（支集）一一对应。
\end{solution}

\begin{solution}{exo:b2:structures:2}
逐点计算，先作用右边的因子。
$\sigma\tau$ 将 $1 \mapsto \sigma(1) = 4$，$\;2 \mapsto \sigma(3)
= 5$, $\;3 \mapsto \sigma(7) = 7$, $\;4 \mapsto \sigma(4) = 2$, $\;5
\mapsto \sigma(5) = 3$, $\;6 \mapsto \sigma(6) = 1$, $\;7 \mapsto
\sigma(2) = 6$:
\[
\sigma\tau = (1\,4\,2\,5\,3\,7\,6),
\]
是一个 $7$-轮换。类似地，$\tau\sigma$ 将 $1 \mapsto \tau(4) = 4$，
$\;2 \mapsto \tau(6) = 6$, $\;3 \mapsto \tau(5) = 5$, $\;4 \mapsto
\tau(2) = 3$, $\;5 \mapsto \tau(3) = 7$, $\;6 \mapsto \tau(1) = 1$,
$\;7 \mapsto \tau(7) = 2$:
\[
\tau\sigma = (1\,4\,3\,5\,7\,2\,6),
\]
同样是 $7$-轮换（符合预期：$\sigma\tau$ 与 $\tau\sigma$ 共轭，
因而有相同的轮换型）。

阶与符号：$\sigma$ 的轮换型为 $(4,2)$，其阶为
$\operatorname{lcm}(4,2) = 4$，符号为 $(-1)^3(-1)^1 = +1$；$\tau$
是 $3$-轮换：阶为 $3$，符号为 $+1$；两个乘积
均为 $7$-轮换：阶为 $7$，符号为 $(-1)^6 = +1$。

$\sigma^{2026}$：$2026 = 4 \times 506 + 2$，故 $\sigma^{2026} =
\sigma^2 = (1\,2)(4\,6)$（对 $4$-轮换取平方，对换
其平方为恒等置换）。
\end{solution}

\begin{solution}{exo:b2:structures:3}
$\{0,1\}^{\N} \to \mathcal{P}(\N)$：数列对应其支集
--- 双射（示性函数），无需定理。

$\{0,1\}^{\N} \to \intcc{0}{1}$：三进制（以 $3$ 为底）映射 $(a_n) \mapsto
\sum 2a_n 3^{-n-1}$ 为单射（两个不同数列首次在
第 $N$ 位不同；尾部无法弥补 $2\cdot
3^{-N-1}$，因为 $\sum_{n > N} 2\cdot 3^{-n-1} = 3^{-N-1} <
2\cdot3^{-N-1}$).

$\intcc{0}{1} \to \{0,1\}^{\N}$：二进展开，例如取
不以全为 $1$ 结尾的那种展开：单射。

由康托–伯恩斯坦定理（\cref{thm:b2:structures:cantorbernstein}）
对最后两个单射应用，$\intcc{0}{1}$ 与
$\{0,1\}^{\N}$ 等势，因而三个集合皆等势。
\end{solution}

\begin{solution}{exo:b2:structures:4}
设 $c = ab = ba$，$d = \operatorname{ord}(c)$。首先 $c^{mn} =
a^{mn} b^{mn} = e$（交换性允许拆分幂），故 $d
\mid mn$。反之，由 $c^d = e$ 得 $a^d = b^{-d}$；该元素
属于 $\langle a\rangle \cap \langle b\rangle$，它是一个子群，其
阶同时整除 $m$ 和 $n$（在各循环群中用拉格朗日定理），
因而是平凡的：$a^d = b^d = e$，故 $m \mid d$ 且 $n \mid d$，
由互素性得 $mn \mid d$。从而 $d = mn$。

在 $\mathfrak{S}_3$ 中，取 $a = (1\,2)$（阶为 $2$）及 $b =
(1\,2\,3)$（阶为 $3$），阶互素但不交换：$ab =
(2\,3)$ 的阶为 $2 \neq 6$ --- 事实上 $\mathfrak{S}_3$ 中没有
阶为 $6$ 的元素。交换性不可或缺。
\end{solution}

\begin{solution}{exo:b2:structures:5}
将每个 $x \in G$ 与 $x^{-1}$ 配对。形如 $\{x, x^{-1}\}$ 的配对中
$x \neq x^{-1}$ 者含两个元素，且划分其并集；其余元素恰为
满足 $x = x^{-1}$ 者，即 $x^2 =
e$ 的元素。由于 $\abs G$ 为偶数，而两元素配对所覆盖的元素个数为偶数，
故集合 $\{x : x^2 = e\}$ 的元素个数为偶数；
它含 $e$，因此至少含另一元素 $x \neq
e$ ——即一个阶为 $2$ 的元素。
\end{solution}

\begin{solution}{exo:b2:structures:6}
$A_n$ 的每个元素都是偶数个
对换之积（\cref{thm:b2:structures:signature}：分解为对换；
对换个数为偶数，因为符号为 $+1$）。只需将
每两个对换之积都用
$3$-轮换表示：
\[
(a\,b)(a\,c) = (a\,c\,b),
\qquad
(a\,b)(c\,d) = (a\,c\,b)(a\,c\,d) \quad (\text{互不相同 } a,b,c,d),
\]
（代入验证），且 $(a\,b)(a\,b) = \mathrm{id}$。因此
$3$-轮换生成 $A_n$。
\end{solution}

\begin{solution}{exo:b2:structures:7}
\emph{$(\Q,+) \to (\Z,+)$：}只有零同态。对任意 $x$ 及
每个 $n \geq 1$，$f(x) = n f\bigl(\frac xn\bigr)$ 是
$n$ 的倍数（在 $\Z$ 中）；能被每个 $n$ 整除的整数只有 $0$，因此
$f(x) = 0$，对所有 $x$ 成立。

\emph{$(\Z/n\Z, +) \to (\Z/m\Z, +)$：}同态由
$c = f(\overline 1)$ 确定，它须满足 $n c \equiv 0 \pmod m$，
即 $c$ 是 $\frac{m}{\gcd(m,n)}$ 的倍数；这样的类共有
$\gcd(m,n)$ 个，且每种选择都确实给出一个同态
（由泛性质，映射 $k \mapsto kc$ 经 $\Z/n\Z$ 分解）。

\emph{$(\Q, +) \to (\Q_+^*, \times)$：}只有平凡同态。若
$f(x) = y$，则对每个 $n$，$y = f(n \cdot \frac xn) =
f(\frac xn)^n$ 是 $n$ 次幂，在 $\Q_+^*$ 中。但有理数
$y \neq 1$ 不能对所有 $n$ 都是 $n$ 次幂：某个素数在
$y$ 中以非零指数 $v$ 出现，则 $n \nmid v$，当 $n > \abs
v$ 时，$n$ 次幂的指数是 $n$ 的倍数（由
唯一分解）。因此 $f \equiv 1$。
\end{solution}

\begin{solution}{exo:b2:structures:8}
$360 = 2^3 \cdot 3^2 \cdot 5$：
$\varphi(360) = 360\bigl(1 - \tfrac12\bigr)\bigl(1 -
\tfrac13\bigr)\bigl(1 - \tfrac15\bigr) = 360 \cdot \tfrac12 \cdot
\tfrac23 \cdot \tfrac45 = 96$。

同余方程组：模数 $8, 9, 5$ 两两互素，总模为 $360$。由 $x
\equiv 3 \pmod 8$ 与 $x \equiv 5 \pmod 9$：设 $x = 3 + 8k$，且 $3 +
8k \equiv 5 \pmod 9$，即 $-k \equiv 2$，$k \equiv -2 \equiv 7
\pmod 9$，于是 $x \equiv 3 + 56 = 59 \pmod{72}$。再令 $59 + 72\ell \equiv
2 \pmod 5$：$4 + 2\ell \equiv 2$，$2\ell \equiv 3 \equiv 8$，$\ell
\equiv 4 \pmod 5$，于是 $x \equiv 59 + 288 = 347 \pmod{360}$。

$3^{2026}$ 的末两位数字：模 $4$ 时，$3^{2026} = 9^{1013} \equiv
1$。模 $25$ 时，$\varphi(25) = 20$，且 $2026 = 20\cdot101 + 6$，故
$3^{2026} \equiv 3^6 = 729 \equiv 4 \pmod{25}$。解同余方程组 $x \equiv 1
\pmod 4$，$x \equiv 4 \pmod{25}$：令 $x = 4 + 25k \equiv 1 \pmod 4$
得 $k \equiv 1 \pmod 4$，于是 $x \equiv 29 \pmod{100}$。所求的末两位
数字为 $29$。
\end{solution}

\begin{solution}{exo:b2:structures:9}
设 $A$ 为有限整环，$a \in A$，$a \neq 0$。映射
$x \mapsto ax$ 是单射（$ax = ay \implies a(x - y) = 0
\implies x = y$，无零因子）；有限集到自身的单射
是满射（见本科一年级分册，鸽巢原理的等价形式）。
于是 $1 = ab$ 对某个 $b$ 成立：每个非零元都可逆，$A$
为域。

$\Z/n\Z$：若 $n$ 为素数，则它是整环（$n \mid ab
\implies n \mid a$ 或 $n \mid b$，欧几里得引理），且有限，故
为域；若 $n = rs$ 为合数，则 $\overline r\,\overline s =
\overline 0$ 给出零因子。
\end{solution}

\begin{solution}{exo:b2:structures:10}
设 $m = \max\{\operatorname{ord}(x) : x \in G\}$，该最大值在 $a$ 处取到。

\emph{断言：每个 $x \in G$ 的阶都整除 $m$。}假设有某个
$x$ 的阶为 $q$，且 $q \nmid m$，则存在素数幂 $p^k$
整除 $q$ 但不整除 $m$。记 $m = p^j m'$，其中 $p \nmid m'$ 且 $j
< k$。元素 $a^{p^j}$ 的阶为 $m'$，元素
$x^{q/p^k}$ 的阶为 $p^k$；这两个阶互素，且二者
可交换（$G \subseteq K^*$ 为阿贝尔群），由
\cref{exo:b2:structures:4}可知它们之积的阶为 $p^k m' > p^j m'
= m$，与最大性矛盾。

因此所有 $x \in G$ 满足 $x^m = 1$：多项式 $X^m - 1$ 至少
有 $\abs G$ 个根，都在域 $K$ 中，从而 $\abs G \leq m$
（非零的 $m$ 次多项式至多有 $m$ 个根，见本科一年级
分册）。但由拉格朗日定理，$m = \operatorname{ord}(a) \leq \abs G$。
因此 $m = \abs G$，而 $\langle a \rangle$ 的势为 $m =
\abs G$，恰为 $G$，从而循环。

对 $K = \Z/p\Z$：$(\Z/p\Z)^*$ 是 $K^*$ 的有限子群，因而
为循环群（阶为 $p - 1$）。
\end{solution}

\begin{solution}{exo:b2:structures:11}
\emph{非循环：}子群 $\langle \frac pq\rangle$ 由
$\frac pq$ 的整数倍组成，它们的
（既约时）分母整除 $q$；因此它不含
$\frac{1}{2q}$。任何单个生成元都无法达到
$\Q$ 中无界的分母。

\emph{非有限生成：}由
$\frac{p_1}{q_1}, \dots, \frac{p_k}{q_k}$ 生成的子群由有理数组成，
其分母整除 $Q = q_1 \cdots q_k$（整数组合的
分母仍整除 $Q$）：它不含 $\frac{1}{2Q}$。

\emph{有限生成的子群是循环的：}记 $Q$ 如上，
子群 $H = \langle \frac{p_1}{q_1}, \dots,
\frac{p_k}{q_k}\rangle$ 含于 $\frac{1}{Q}\Z$。映射 $x
\mapsto Qx$ 是从 $\frac1Q\Z$ 到 $\Z$ 的同构，它把
$H$ 映为 $\Z$ 的子群，它形如 $n\Z$，$n$ 为某个整数（见本科一年级
分册）。于是 $H = \frac{n}{Q}\Z$ 为循环群，由
$\frac nQ$ 生成。
\end{solution}

\begin{solution}{exo:b2:structures:12}
\emph{存在可数子集。}设 $E$ 为无限集，归纳构造 $a_0,
a_1, a_2, \dots$：$E$ 非空，取 $a_0 \in E$；
若 $a_0, \dots, a_n$ 已取定，则 $E \setminus \{a_0, \dots,
a_n\}$ 非空（$E$ 非有限），在其中取 $a_{n+1}$。
由构造，$a_n$ 两两不同，因此 $A = \{a_n : n
\in \N\}$ 为 $E$ 的可数子集。

\emph{无限集 $\implies$ 与某个真子集等势。}定义
映射 $f \colon E \to E \setminus \{a_0\}$，$f(a_n) = a_{n+1}$，且
$f(x) = x$（$x \notin A$）。它为单射（两段都是单射，
且像不相交），且满射到 $E \setminus
\{a_0\}$：每个 $a_{n+1}$ 与每个 $x \notin A$ 都在像集中。因此
$E$ 与真子集 $E \setminus \{a_0\}$ 等势。

\emph{逆命题。}若 $E$ 为有限集，且 $g \colon E \to F$ 为
到 $F \subseteq E$ 上的双射（$F \neq E$），则 $g$ 为 $E$ 到自身的
不满射的单射，
与鸽笼原理相矛盾（本科一年级分册：有限集的
单射自映射必为双射）。因此，与真子集等势的集合
为无限集。
\end{solution}

\begin{solution}{pb:b2:structures:1}
\textbf{1.} 构形为对 $16$ 个格子各分配
$16$ 种内容之一（数码块 $1$--$15$ 或空格 $b = 16$），
且每种内容恰用一次，即为双射 $\intint1{16} \to
\intint1{16}$，是 $\mathfrak{S}_{16}$ 中的元素；共有 $16!
= 20\,922\,789\,888\,000$ 个。合法移动是将一块数码块滑到空格旁，
故移动次数即空格所在格子的邻格数：
四个角格的邻格数为 $2$，
八个边格为 $3$，四个内部格为 $4$。

\textbf{2.} 滑动之后，格子 $p$ 中为原先格子
$c$ 中的内容，格子 $c$ 中为空格，其余格子保持不变：
$\sigma'(p) = \sigma(c)$，$\sigma'(c) = \sigma(p) = 16$，
$\sigma' = \sigma$ 在其余位置成立，即 $\sigma' = \sigma
\circ (p\ c)$。由于 $\varepsilon$ 为同态，且
$\varepsilon\bigl((p\ c)\bigr) = -1$：
$\varepsilon(\sigma') = -\varepsilon(\sigma)$。

\textbf{3.} 相邻格子恰在两个坐标之一上相差一步，
因此 $i + j$ 的奇偶性改变，$\chi$
在相邻格子上取相反的值。一次移动将空格
从 $p$ 移到相邻的 $c$，使 $\chi(\text{空格所在格})$ 翻转。
空格沿闭合路径移动时，$\chi$ 每次移动翻转一次，
并回到初始值，故移动次数为偶数。

\textbf{4.} 由第 2、3 问可知，一次移动同时翻转
$I(\sigma) = \varepsilon(\sigma)\chi(\sigma^{-1}(16))$ 的两个因子；其
乘积不变。对于已复原的构形：
$\varepsilon(\mathrm{id}) = +1$，且空格位于格子 $16$，
第 $4$ 行第 $4$ 列：$\chi(16) = (-1)^{8} = +1$，故
$I(\mathrm{id}) = +1$。

\textbf{5.} $\sigma_L$ 是格子间的对换 $(14\ 15)$：
$\varepsilon(\sigma_L) = -1$；其空格在原位，$\chi(16) = +1$：
$I(\sigma_L) = -1 \neq +1 = I(\mathrm{id})$。由于 $I$ 在每次移动下
保持不变，故没有任何移动序列能使 $\sigma_L$
与 $\mathrm{id}$ 相连。奖品在结构上是安全的。

\textbf{6.} 固定格子 $p$，取两个格子 $c \neq d$，二者均异于 $p$，
并令 $\tau_0 = (c\ d)$。在空格位于 $p$ 的构形集合上，
映射 $\sigma \mapsto \sigma
\circ \tau_0$ 为对合（它保持 $\sigma(p) = 16$，
因 $\tau_0$ 固定 $p$）且使 $\varepsilon$ 变号，从而使 $I$ 变号，
它将 $I = +1$ 的构形与
$I = -1$ 的构形双射配对。因此 $16$ 个空格位置中的每一个
各贡献 $15!/2$ 个 $I = +1$ 的构形，于是
\[
\abs{\{I = +1\}} = 16 \cdot \frac{15!}{2} = \frac{16!}{2}.
\]

\textbf{7.} 将 $c$ 处数码块滑入 $p$ 的移动，
可由将同一数码块（现位于 $p$）滑回 $c$ 而撤销；复合
$(p\ c)$ 两次即为恒等映射。由此得：自反性（空序列）、
对称性（将序列逆置，逐一撤销每次移动）、
传递性（将序列接连）：为等价关系。每个
$\sigma \in R$ 有 $I(\sigma) = I(\mathrm{id}) = +1$，由第 
4 问，故 $R \subseteq \{I = +1\}$；又 $\sigma_L \notin R$，给出
第二个等价类。

\textbf{8.} 若 $\sigma(16) = 16$，则 $\sigma$ 置换
格子 $1, \dots, 15$，记此限制为 $\rho$。在
末尾添加一个不动点，既不改变轮换型，也不改变符号
（将 $\rho$ 分解为对换；同一乘积在
$\mathfrak{S}_{16}$ 中仍有效），故 $\varepsilon(\sigma) =
\varepsilon(\rho)$，又 $\chi(16) = +1$，得 $I(\sigma) =
\varepsilon(\rho)$。任一构形都可化为
空格在原位的构形：方格图连通，故沿相邻格子的路径
将空格移到格子 $16$（每一步均为合法移动）。
现设每个偶置换 $\rho \in \mathfrak{S}_{15}$ 都
可由某个程序实现。给定 $\sigma$，且 $I(\sigma) = +1$：将空格走回
原位，得到 $\widetilde\sigma$（与 $\sigma$ 等价），
其中 $I(\widetilde\sigma) = +1$，即其限制 $\rho$ 为
偶置换；实现 $\rho$ 的程序将 $\widetilde\sigma$ 变为
$\widetilde\sigma \circ \rho^{-1} = \mathrm{id}$（见第
9 问）。由传递性得 $\sigma \in R$，从而 $\{I = +1\} \subseteq
R$，两者相等。

\textbf{9.} \emph{单步移动：}$c$ 中的内容移至 $p$，
空格移至 $c$：效果为 $\pi = (p\ c)$，且确有
$\sigma' = \sigma \circ (p\ c) = \sigma \circ \pi^{-1}$。
\emph{归纳：}若一序列的效果为 $\pi_1$，并将
$\sigma$ 变为 $\sigma \circ \pi_1^{-1}$，在其后接一次
效果为 $\pi_2 = (p'\ c')$ 的移动，得到 $(\sigma \circ \pi_1^{-1})
\circ \pi_2^{-1} = \sigma \circ (\pi_2\pi_1)^{-1}$，内容按
$\pi_2 \circ \pi_1$ 移动（先 $\pi_1$，后 $\pi_2$）。因此
效果可复合，从 $\sigma$ 出发运行程序，终止于 $\sigma
\circ \pi^{-1}$。\emph{子群：}空程序的效果为
$\mathrm{id}$；接连得到乘积；将程序逆置
（第 7 问）得到逆元。程序的效果固定格子 $16$
（空格始末均在原位），故 $H \leq \mathfrak{S}_{15}$。
\emph{偶性：}含 $k$ 次移动的程序，其 $k$ 为偶数（第
3 问），且 $\varepsilon(\sigma \circ \pi^{-1}) =
(-1)^k\varepsilon(\sigma)$ 迫使 $\varepsilon(\pi) = +1$，即 $H
\subseteq A_{15}$。

\textbf{10.} 从位于 $16$ 的空格出发，追踪四次滑动：
移动 $16 \to 12$ 将 $12$ 中的内容送至 $16$；移动 $12 \to 11$
将 $11$ 中的内容送至 $12$；移动 $11 \to 15$ 将
$15$ 中的内容送至 $11$；移动 $15 \to 16$ 将停放在
$16$ 中的内容（原位于 $12$）送至 $15$。总效果为 $11 \mapsto
12$，$12 \mapsto 15$，$15 \mapsto 11$，空格回到原位：效果为
$(11\ 12\ 15)$。逆向走完全程可将其抵消：效果为 $(11\ 12\
15)^{-1} = (11\ 15\ 12)$。两者均为程序的效果，故
属于 $H$。

\textbf{11.} 相邻格的关系：在所列每一数对中，
两格在同一行内相差 $1$（$16{-}15$，
$15{-}14$，$14{-}13$；$1{-}2$，$2{-}3$，$3{-}4$；$8{-}7$，
$7{-}6$；$10{-}11$，$11{-}12$），或在同一列内相差 $4$
（$13{-}9$，$9{-}5$，$5{-}1$；$4{-}8$；$6{-}10$；$12{-}16$）：这是一条
经过全部 $16$ 个格的闭路径，长度为 $16$。效果与
第 10 题相同，记所经过的格依次为 $c_0 = 16, c_1 = 15,
\dots, c_{15} = 12$：$c_i$ 中的内容移至 $c_{i-1}$，
对 $i = 2, \dots, 15$，而 $c_1$ 中的内容在第一次移动后停放于 $16$ 中，
由最后一次移动送至 $c_{15}$。
于是效果映射为 $15 \mapsto 12$，$14 \mapsto 15$，$13
\mapsto 14$，$9 \mapsto 13$，$5 \mapsto 9$，$1 \mapsto 5$，$2
\mapsto 1$，$3 \mapsto 2$，$4 \mapsto 3$，$8 \mapsto 4$，$7
\mapsto 8$，$6 \mapsto 7$，$10 \mapsto 6$，$11 \mapsto 10$，$12
\mapsto 11$：恰为 $15$-轮换 $\zeta$。其轮换次序
始于 $x_0 = 15$，$x_1 = 12$，$x_2 = 11$，而 $(x_0\ x_1\ x_2)
= (15\ 12\ 11)$ 给出 $15 \mapsto 12 \mapsto 11 \mapsto 15$ ---
恰为 $(11\ 15\ 12)$，即逆向的基本巡回。

\textbf{12.} 设 $\gamma = (a\ b\ c)$ 及 $x \in \intint1n$。若
$x = g(a)$ 时：$g\gamma g^{-1}(x) = g(\gamma(a)) = g(b)$；同理
$g(b) \mapsto g(c)$，$g(c) \mapsto g(a)$。若 $x \notin
\{g(a), g(b), g(c)\}$，则 $g^{-1}(x) \notin \{a,b,c\}$
被 $\gamma$ 固定，故 $x$ 亦被固定。因此 $g\gamma g^{-1} =
(g(a)\ g(b)\ g(c))$。又对 $g, h \in H$，$ghg^{-1} \in H$
由子群公理可知。

\textbf{13.} $\zeta \in H$（第 11 题）及 $s_0 = (x_0\ x_1\
x_2) \in H$（第 10--11 题）。由 $\zeta(x_i) = x_{i+1}$
（指标按模 $15$ 理解），由第 12 题得
\[
\zeta^{t}\,s_0\,\zeta^{-t}
= \bigl(\zeta^t(x_0)\ \zeta^t(x_1)\ \zeta^t(x_2)\bigr)
= (x_t\ x_{t+1}\ x_{t+2}) = s_t \in H
\qquad (t = 0, 1, \dots, 14).
\]

\textbf{14.} 不计求逆，可设 $s = (a\ b\ c)$ 及 $t =
(b\ c\ d)$（$3$-轮换在 $\{a,b,c\}$ 上为 $(a\ b\ c)$ 或其
逆；$\{b,c,d\}$ 上同理；以其逆替换生成元，
不改变 $\langle s, t\rangle$）。于是
先作用 $t$，则
\[
st \colon a \mapsto b,\quad b \mapsto a,\quad c \mapsto d,\quad
d \mapsto c, \qquad\text{即}\quad st = (a\ b)(c\ d),
\]
这是一个双对换。子群 $G = \langle s, t\rangle$
由四个字母的偶置换构成，故 $G \leq
A_4$，且 $\abs G \mid 12$；它含有阶为 $3$ 的元素
及阶为 $2$ 的元素，故 $6 \mid \abs G$（拉格朗日，
\cref{thm:b2:structures:lagrange}，应用于两个循环
子群）。若 $A_4$ 有子群 $K$，其阶为 $6$，则其
指数为 $2$，于是 $g^2 \in K$ 对一切 $g \in A_4$ 成立：对
$g \in K$ 显然；对 $g \notin K$，陪集仅有
$K$ 与 $gK$，故陪集 $g^2K$ 为 $K$ 或 $gK$，而 $g^2K =
gK$ 将迫使 $g \in K$。故每个平方皆在 $K$ 中。但每个
$3$-轮换 $\gamma$ 皆为平方，$\gamma = (\gamma^2)^2$，而
$A_4$ 含有八个 $3$-轮换：$8 > 6$，矛盾。因此
$\abs G = 12$，即 $G = A_4$。

\textbf{15.} 将 $u \mapsto a$，$v \mapsto b$ 延拓为一个双射
$g_0$（$X$ 上的置换，将其余 $k - 2$ 个字母双射地
任意映到 $\{a, b\}$ 的补集上）。若 $g_0$ 为奇置换，
任取两个不同字母 $s_1, t_1 \in X \setminus \{u, v\}$
（可行，因 $k \geq 4$），将 $g_0$ 换为 $g_0 \circ (s_1\
t_1)$，它是偶置换，且仍有 $u \mapsto a$，$v \mapsto
b$。在 $X$ 之外取恒等映射延拓，得偶置换 $g \in
G$（它是 $X$ 上的偶置换）。由第 12 题得：
\[
g\,(u\ v\ w)\,g^{-1} = (g(u)\ g(v)\ g(w)) = (a\ b\ w) \in G,
\]
利用 $g(w) = w$。

\textbf{16.} 每个 $3$-轮换（支撑在 $X \cup \{w\}$ 中）皆属于 $G$：
支撑在 $X$ 中者都是 $X$ 上的偶置换；其中支撑为
$\{a, b, w\}$ 者为 $(a\ b\ w)$ 或 $(b\ a\ w)$，二者皆由
第 15 题给出。由\cref{exo:b2:structures:6}，
$3$-轮换生成 $(k+1)$ 元集 $X \cup \{w\}$
的交错群，故 $G$ 包含
$X \cup \{w\}$ 的全部偶置换。\emph{逐步推进：}令 $G = \langle s_0, \dots,
s_{12}\rangle$。对 $s_0 = (x_0\ x_1\ x_2)$ 及
$s_1 = (x_1\ x_2\ x_3)$（二者支撑共享 $\{x_1, x_2\}$），应用引理 A 得
$X_4 = \{x_0, x_1, x_2, x_3\}$ 的全部偶置换。若 $G$
包含 $X_m = \{x_0, \dots,
x_{m-1}\}$ 的全部偶置换（$4 \leq m \leq 14$），则 $s_{m-2} = (x_{m-2}\
x_{m-1}\ x_m)$ 中 $u = x_{m-2}, v = x_{m-1} \in X_m$，并引入新字母
$w = x_m$：由引理 B 及第一部分，得
$X_{m+1}$ 的全部偶置换。归纳至 $m = 14$，得 $G \supseteq
A_{15}$（即全部十五个格的偶置换），又 $G
\subseteq A_{15}$，因每个 $s_t$ 皆为偶置换：$\langle s_0, \dots,
s_{12}\rangle = A_{15}$。

\textbf{17.} 由第 13 题与第 16 题：$A_{15} = \langle s_0, \dots,
s_{12}\rangle \subseteq H$；由第 9 题，$H \subseteq A_{15}$。故
$H = A_{15}$，其阶为 $15!/2 = 653\,837\,184\,000$：
十五块方块的每个偶重排都是某个程序的效果。

\textbf{18.} 第 8 题将 $R = \{I = +1\}$ 归结为
对每个偶 $\rho \in \mathfrak{S}_{15}$ 以程序实现：
由第 17 题完成。由第 6 题，$\abs R = 16!/2 =
10\,461\,394\,944\,000$。\emph{两类：}令 $t_0 = (14\
15)$ 作用于\emph{内容}：$\varphi(\sigma) = t_0 \circ
\sigma$。自 $\sigma$ 出发的合法移动，亦是
$\varphi(\sigma)$ 出发的合法移动（空格所在格不变：
$(t_0\sigma)^{-1}(16) = \sigma^{-1}(t_0(16)) = \sigma^{-1}(16)$，
被移动的格亦同），且 $\varphi(\sigma \circ \tau)
= \varphi(\sigma) \circ \tau$：$\varphi$ 将移动序列映为
移动序列，且为双射（它是对合）。它使 $I$ 变号：
$\varepsilon(t_0\sigma) = -\varepsilon(\sigma)$，空格所在格
相同。因此 $\varphi$ 将类 $R = \{I = +1\}$
（$\mathrm{id}$ 所在的类）双射地映到 $\varphi(\mathrm{id})
= \sigma_L$ 所在的类，故其即为全部 $\{I = -1\}$，于是恰有两
类。这就是约翰逊--斯托里定理。

\textbf{19.} 按阅读顺序为各格编号，令 $k = 4(i -
1) + j$ 为空格所在格。数 $\sigma$ 的逆序
（即满足 $x < y$ 且 $\sigma(x) > \sigma(y)$ 的格对）：两格皆含牌的格对
各贡献 $N$；涉及空格的格对中，空格之后的各格
所含之牌皆 $< 16$，各构成逆序（$16 - k$
对），空格之前的格从不构成逆序。故
$\varepsilon(\sigma) = (-1)^{N + 16 - k} = (-1)^{N + k}$。
由于 $k = 4(i-1) + j \equiv j \pmod 2$，
\[
I(\sigma) = (-1)^{N + j}\,(-1)^{i + j} = (-1)^{N + i}
= (-1)^{N + r + 1}
\]
利用 $i = 5 - r$。由第 18 题，$\sigma$ 可解当且仅当
$I(\sigma) = +1$ 当且仅当 $N + r$ 为奇数。验算：已解局面 $N = 0$，$r
= 1$，为奇数，可解；洛伊德局面 $N = 1$，$r = 1$，为偶数，不可解。

\textbf{20.} \emph{作用：}$e \cdot \sigma = \sigma \circ
\mathrm{id} = \sigma$，且 $g \cdot (h \cdot \sigma) = \sigma
\circ h^{-1} \circ g^{-1} = \sigma \circ (gh)^{-1} = (gh) \cdot
\sigma$，又 $\sigma \circ h^{-1}$ 仍为空格归位的
局面（$h$ 固定格 $16$）。\emph{自由性：}$\sigma \circ
h^{-1} = \sigma$ 推出 $h^{-1} = \mathrm{id}$（右乘
$\sigma^{-1}$）。\emph{轨道即程序类：}由第 9 题
指出，由 $\sigma$ 经程序可达的局面恰为
$\sigma \circ \pi^{-1}$，$\pi \in H$，即轨道 $H
\cdot \sigma$。\emph{计数：}由自由性，$h \mapsto h \cdot
\sigma$ 为单射，故每条轨道含 $\abs H = 15!/2$ 个元素；
于是 $15!$ 个空格归位的局面被分成
$15!\,/\,(15!/2) = 2$ 条轨道 --- 即
两个约翰逊--斯托里类在空格归位情形下的投影。

\textbf{21.} 对棋盘染色，$3 \times 3$ 方格图是二部图：
即行走的每一步都改变颜色，故
每条\emph{闭}行走长度为偶数。若闭行走访问
$9$ 个格各恰一次，则长度为 $9$，为奇数：
不可能。故第三部分的大环游构造在
八数码问题上不可用。

\textbf{22.} \emph{周边巡回} $9 \to 8 \to 7 \to 4 \to 1
\to 2 \to 3 \to 6 \to 9$（每步皆相邻，长度 $8$，为偶数）：
由第 11 题的记账法，取 $c_1 = 8, c_2 = 7, c_3 =
4, c_4 = 1, c_5 = 2, c_6 = 3, c_7 = 6$，效果为
\[
\zeta' = (8\ 6\ 3\ 2\ 1\ 4\ 7),
\]
一个 $7$-轮换，固定中心 $5$（$7$ 的牌移到 $8$，
$4$ 的牌移到 $7$，$1$ 的牌移到 $4$，$2$ 的牌移到 $1$，$3$ 的牌移到 $2$，
$6$ 的牌移到 $3$，$8$ 的牌移到 $6$）。\emph{角点巡回} $9 \to 6 \to
5 \to 8 \to 9$：效果为 $(6\ 8\ 5)$（$5$ 的牌移到
$6$，$8$ 的牌移到 $5$，$6$ 的牌（停放于 $9$）移到 $8$）。
令 $y_t = \zeta'^{\,t}(8)$：$y_0 = 8, y_1 = 6, y_2 = 3, y_3 = 2,
y_4 = 1, y_5 = 4, y_6 = 7$。共轭（第 12 题）：
\[
\zeta'^{\,t}\,(6\ 8\ 5)\,\zeta'^{-t}
= (y_{t+1}\ y_t\ 5) =: T_t \in H_{3\times3},
\]
由于 $\zeta'$ 固定 $5$，$T_0 = (y_1\ y_0\ 5)$ 的支集与
$T_1 = (y_2\ y_1\ 5)$ 的支集恰共有 $\{y_1, 5\}$：由引理 A
给出 $\{y_0, y_1, y_2, 5\}$ 上的全部偶置换。随后
$T_2 = (y_3\ y_2\ 5)$ 由引理 B 加入 $y_3$（其字母
$y_2, 5$ 位于当前集合中，$k = 4$），$T_3, T_4, T_5$
依次加入 $y_4, y_5, y_6$：八个非归位格上的全部
偶置换皆属于程序群，程序群亦只含
偶置换（第 9 题的论证与棋盘无关）。
故 $H_{3\times3} = A_8$，又
第 6、8、18 题的论证（亦与棋盘无关）表明，可达局面
恰为满足 $I = +1$ 的那些，即 $9!$ 的一半，
共 $181\,440$ 个。

\textbf{23.} 将各格沿圈依次编号为 $0, \dots, n-1$。
一步移动即将空格与其两个相邻格之一交换。自空格之后起按圈序读取各牌，
记为字 $w$：
列出 $n - 1$ 块牌。空格前移一步，
相当于将 $(p, w)$ 换为 $(p + 1, \rho w)$，其中 $p$ 为空格所在格，
$\rho$ 将字循环地旋转一位；后退一步
则为其逆。于是牌的\emph{循环}次序（即字在旋转下的类）
是不变量。$(p, w)$ 所在的可达类
恰为映射 $g \colon (p, w) \mapsto (p+1,
\rho w)$ 的轨道，即两个循环群之积中阶为 $\operatorname{lcm}(n, n-1) =
n(n-1)$ 的元素（$\Z/n\Z$ 的平移
与 $n-1$ 个字位的旋转），其中
最小公倍数为 $n(n-1)$，因 $\gcd(n, n-1) = 1$：每一类
恰含 $n(n-1)$ 个局面，且均具有同一项链。
类数为 $n!\,/\,\bigl(n(n-1)\bigr) = (n-2)!$。当 $n \geq 5$ 时，
$(n-2)! > 2$，奇偶性不变量（至多两类）对障碍的
绝大部分视而不见；$4
\times 4$ 棋盘的丰富性——在那里奇偶性是\emph{唯一}障碍——
乃真正几何性的事实，而非形式上的。

\textbf{24.} 两个托盘中的牌均为完全逆序，
故两种情形下 $N = \binom{15}{2} = 105$（每一对牌皆构成逆序），
\emph{空格归位：}$r = 1$，$N + r = 106$ 为偶数，
不可解。\emph{空格位于格 $1$：}空格位于顶行，
$r = 4$，$N + r = 109$ 为奇数：可解。两个仅因空洞位置而异的托盘
分处墙的两侧。

\textbf{25.} \emph{同态性质：}它把“一步移动 =
一次对换”化为“一步移动 = 一次符号翻转”（第
2、4 题），使 $I$ 可逐步计算。\emph{拉格朗日：}它
在引理 A 中迫使 $6 \mid \abs{\langle s, t\rangle}$，并
确定了阶为 $6$ 的排除中各陪集的大小（第 14 题）。
\emph{$3$-轮换生成：}它把“$H$ 含有
足够多的 $3$-轮换”化为“$H$ 含有全部 $A_{15}$”
（第 16 题）。\emph{共轭：}它造出十五个连续的
$3$-轮换，其源自单个 $2 \times 2$ 巡回，
经大环游搬运（第 12--13 题），以及引理 B 中的
$3$-轮换 $(a\ b\ w)$。\emph{元原理：}一个
不变量证明不可能，显式构造证明
可能；一个问题恰在两界相遇时才算完全解决，
两界在此相遇于一半。
\end{solution}
